\scanpage{111}
\uline{Corollaire} Considérons le sous grpe \inserted{(invariant) $H^{\unsure{0|\circ}}$} du grpe \deleted{$\unsure{G}$} \inserted{$H={}$}$\underline{\mathrm{Aut}}(S,A\,;a,b,c)$ \deleted{engendré} \inserted{formé} par les produits d'un nb pair de symétries $\sigma_{\underline t}$ relatifs à des triangles verticaux.

Alors $H^0$ est transitif sur l'ens $\widetilde T$ des triangles supérieurs. De façon précise, si $\deleted{\unsure{\underline s,\underline t'}}\inserted{\underline s,\underline t}\in\widetilde T$, et si $a\in\underline t_0$, soit \deleted{$u\in$} $b\in t_0$, $b\neq a$ tel que $\exists\,u\in\widetilde E_b$ lié aux deux sommets \inserted{$s,t$} de $\underline s\ \underline t$ sur $a$, alors $\exists$

\deleted{produit $U$ de deux symétries \deleted{relatifs} $\sigma_{\unsure{v}},\sigma_{w}$ ($v,w\in E_c=\widetilde E_c/2$) tel que $Us=t$, \inserted{$U\underline s$ et $\underline t$ coïncident \uline{en $a$} ;} et on a ensuite \uline{le choix} pour trouver $V$, produit de deux symétries \unsure{latérales|relatives} en $b$, ou en $c$, tels que $VU\underline s=\underline t$.}

\deleted{Si on le prend en $\underline c$, comme $V$ et $U$ commutent, $VU$ sera encore un autom. involutif}

% archive sheet number: 110
